\documentclass[a4paper, 12pt]{article}
\usepackage[utf8]{inputenc}
\usepackage[english, russian]{babel}
\usepackage{eulervm}

\title{Работа от некого Р}
\author{\textit{Samsyka} R}
\date{March 1, 2025} 

\begin{document}
\maketitle 
\TeX \--- это компьютерная программа, созданная Дональдом Кнутом (Donald E. Knuth). Она предназначена для вёрстки текста и математических формул. Кнут начал писать \TeX в 1977 году из-за расстройства от того, что Американское Математическое Сообщество делало с его статьями в процессе их публикации. Где-то в 1974 году он даже прекратил посылать статьи: \textbf{"<Просто мне было слишком больно смотреть на конечный результат">.} \TeX, в том виде, в котором мы его используем, был выпущен в 1982 году и слегка улучшен с годами. Последние несколько лет \TeX стал чрезвычайно стабилен. Кнут утверждает, что в нем практически нет ошибок. Номер версии \TeX сходится к $\pi$ и сейчас равен 3.14159. \TeX произносится как "<\TeX">.
\end{document}
\begin{equation}
\int \limits_S \left( \frac{\partial Q}{\partial x} - \frac{\partial P}{\partial y} \right)\, dx \, dy =\oint \limits_C P\,dx + Q \, dy
\end{equation}